\chapter{뉴런은 한 가지 소자가 아니다}
% full version: chapters/18_eetypes.tex

1장에서는 뉴런을 어떤 물질을 내보내는지에 따라 분류했다. 엔지니어라면 다르게 분류할 것이다. 각 뉴런이 어떤 소자로 동작하는가. 그러면 똑같아 보이는 겉모습 아래 전자 부품 상자 하나가 통째로 숨어 있음이 드러난다.

\section*{전자 부품 상자}

2장의 구멍 난 양동이는 누설 적분기다. 짧은 시간 창 안의 신호를 더한다. 창이 아주 짧으면 동시성 검출기가 된다. 변화에만 반응하고 곧 조용해지는 뉴런은 변화를 통과시키는 필터다. 스스로 동작하는 뉴런은 메트로놈, 곧 발진기다. 입력이 그 박자를 조정한다면 전압 제어 발진기다. 딸깍을 버스트로 내는 뉴런은 “선”을 쓰는 버스트 발생기다. 음파에 맞춰 딸깍하는 뉴런은 위상 검출기다.

\pencilfig{18}{\pencilart{18}}{똑같은 겉모습 아래 전자 부품 상자 하나가 숨어 있다. 저항, 커패시터, 코일, 발진기.}

지금까지 다루지 않은 소자도 있다. \emph{공진기}는 그네와 비슷하다. 정해진 리듬으로 밀 때 가장 잘 반응한다. 제때 밀어 줘야만 높이 올라가는 그네처럼. 뉴런 안에서는 변화에 저항하는 느린 전류가 용수철 역할을 한다. \emph{누설 없는 적분기}는 누설 없는 커패시터처럼 값을 오래 붙잡아 둔다. 여러분이 옆을 볼 때 눈의 위치를 유지하는 네트워크가 이렇게 일한다. 다만 그러려면 되먹임이 거의 완벽하게 맞춰져 있어야 하고, 이것은 비싼 튜닝이다. \emph{래치}는 짧은 자극 하나를 받아, 누가 끌 때까지 오래 동작하는 뉴런이다. 누르면 고정되는 스위치와 같다. 근육을 제어하는 뉴런에서는 세로토닌이 이 모드를 켠다.

\section*{카멜레온 소자}

이 분류의 가장 큰 발견은 이것들이 서로 다른 부품이 아니라 같은 부품의 서로 다른 \emph{모드}라는 점이다. 같은 뉴런이 낮에는 적분기이고 밤에는 버스트 발생기일 수 있다. 방송국이 무엇을 말하느냐에 달렸다. 엔지니어라면 뇌를 재프로그래밍 가능한 논리 배열이라고 부를 것이다. 회로는 하나이고, 부품의 역할은 설정이 바꾼다.

뉴런이 이 모든 모드로 동작할 수 있다는 것은 측정되었다. 뇌가 모드 전환을 계산에 어떻게 쓰는지는 대부분 가설이다.

\fullversion{18}
